\subsection{I-28. Nonzero endpoint constants fit inside the Grassmannian support budget}
\label{I-28}

\paragraph{Statement.}
The two extra endpoint positions in I-25 are unnecessary when
the endpoint vectors are nonzero at the center.  Let
$1\le r\le w-1$, put $g=r(w-r)$, and consider
\[
 q(Y_0,\ldots,Y_{m+1})
 =\left\langle u(Y_0)
       \prod_{i=1}^m\bigwedge^r A_i(Y_i),\ v(Y_{m+1})\right\rangle.
 \tag{I.39}
\]
Here $A_i(0)$ is invertible, $u(Y_0)$ is a polynomial
decomposable exterior row with $u(0)\ne0$, and $v(Y_{m+1})$
is a polynomial exterior column with $v(0)\ne0$.
Only the left row is required to be decomposable.
All variables displayed are distinct; constant endpoints may
be included by omitting their variables.
If $q\ne0$, setting all but at most $g$ variable positions
to zero leaves a nonzero polynomial.  In particular
$\operatorname{SV}_{g}$ hits this class, and the count includes
both endpoint variables.

\paragraph{Proof.}
Extend the coefficient field to an algebraic closure; a zero
restriction over the original field stays zero after this
extension, so proving existence there suffices.  Use constant
changes of basis to make all $A_i(0)=I_w$, as in I-25.
They preserve nonzero endpoint constants and decomposability.
Restrict each internal parameter to the nonempty open set
where $A_i$ is invertible, and the left parameter to the open
set where $u$ is nonzero.  These opens contain zero.

Let $Z_{-1}$ be the projective point $[u(0)]$, and let $Z_0$
be the closure of the image of $[u(Y_0)]$.  For $i\ge1$,
let $Z_i$ be the closure of the projective image after applying
the first $i$ internal matrices.  Each $Z_i$ is an irreducible
closed subvariety of $\operatorname{Gr}(r,w)$, and
\[
 Z_{-1}\subseteq Z_0\subseteq\cdots\subseteq Z_m.
\]
For the first inclusion, use the available parameter value
$Y_0=0$; for the others, use $A_i(0)=I_w$.

Retain a parameter position exactly when this inclusion is
strict.  There are at most $\dim Z_m$ retained positions,
because a strict inclusion of irreducible closed varieties
strictly increases dimension, starting at a point.
Setting all other positions to zero preserves the final
image closure $Z_m$.  To check this rigorously by induction,
at a retained position the previous restricted image is
dense in the same $Z_{i-1}$; applying the regular invertible
actions and taking closure gives the same $Z_i$.  At a
nonretained position, zero substitution applies the identity
and leaves $Z_{i-1}=Z_i$.  The initial left curve obeys the
same statement: if its image closure is a point, setting its
parameter to zero preserves that point.

Write $\ell_0$ for the nonzero linear functional determined
by $v(0)$.  If $\ell_0$ does not vanish identically on $Z_m$,
set the right parameter to zero as well.  The retained image
is still dense in $Z_m$, so the resulting scalar polynomial
is nonzero, using at most $\dim Z_m\le g$ positions.

Otherwise $Z_m$ lies in the hyperplane $\ell_0=0$.
The Grassmannian is not contained in that hyperplane:
the coordinate decomposable wedges span the full exterior
space, and $\ell_0\ne0$.  Therefore
\[
 \dim Z_m\le \dim\bigl(\operatorname{Gr}(r,w)
                                  \cap\{\ell_0=0\}\bigr)\le g-1.
\]
As $q\ne0$, some value of the right parameter gives a
functional that does not vanish identically on $Z_m$.
The retained prefix image has the same closure, so keeping
this right parameter free gives a nonzero scalar polynomial
on at most $\dim Z_m+1\le g$ positions.  This proves the
restriction assertion in both cases.  A nonzero restriction
to at most $g$ coordinate positions is a specialization of
$\operatorname{SV}_g$, by setting its selectors to the
corresponding Lagrange nodes.  Thus the generator assertion
follows in every characteristic, with the usual distinct-node
field or extension convention.  No algorithm to find the
retained positions from the unknown matrices is assumed.
\hfill$\square$

\paragraph{Application to the width-three construction.}
In the rank-two stage of I-21, a selected nonzero minor of
a local rank test is a contraction of a second-compound
word with an invertible interior.  Its compound row and
column at either endpoint are nonzero at the initial common
center $x_i=t$.  Indeed, a nonzero selected contraction has
nonzero endpoint rows and columns as univariate polynomial
vectors; evaluating a nonzero univariate entry at the
transcendental element $t$ cannot make it zero.  Constant
endpoints are covered directly.  Single-factor tests are
already preserved by the same argument and need no pairs.

Consequently I-28 applies with $w=3,r=2$, giving support
budget $g=2$ for the entire rank-two test, including its
endpoints.  Two fresh SV pairs simultaneously preserve all
rank-two tests: each selected minor is a nonzero polynomial
after substitution, and the product of the finitely many
selected minors is nonzero.  Previously preserved rank-three
minors remain nonzero upon setting the new amplitudes to zero.
The rank-one stage of I-21 still uses its independently
proved six-pair budget.  With disjoint fresh parameter pairs
the two stages therefore give
\[
 \mathcal G_i^{(8)}
      =t+\sum_{j=1}^{8}z_j\lambda_i(y_j),
 \qquad f\ne0\Longrightarrow f(\mathcal G^{(8)})\ne0
 \tag{I.40}
\]
for every width-at-most-three one-hidden-order ROABP in the
scope of I-21.

All degree and interpolation arguments of I-21 remain
unchanged except for replacing ten pairs by eight.  In its
notation, this yields the explicit sample bound
\[
 (D+1)\bigl((n-1)D+1\bigr)^8\binom{D+8}{8},
 \tag{I.41}
\]
over a working field containing at least
$\max\{n,D+1,(n-1)D+1\}$ distinct available elements.
The earlier ten-pair theorem remains valid; (I.40) is a
strictly smaller proved sampler.  This refinement does not
address the singular reset profiles at variable width and
does not claim a polynomial dependence on that width.
